% !TeX root = ../../main.tex

% Configuración para mostrar código Python
\lstset{
    language=Python,
    basicstyle=\ttfamily\small,
    keywordstyle=\bfseries,
    commentstyle=\itshape,
    numbers=left,
    numberstyle=\tiny,
    stepnumber=1,
    numbersep=8pt,
    frame=single,
    breaklines=true,
    showstringspaces=false,
    columns=fullflexible,
    inputencoding=utf8,
    literate=
        {á}{{\'a}}1
        {é}{{\'e}}1
        {í}{{\'i}}1
        {ó}{{\'o}}1
        {ú}{{\'u}}1
        {ü}{{\"u}}1
        {ñ}{{\~n}}1
}

\section{Introducción}

MapReduce es un modelo de programación que permite procesar grandes
cantidades de información dividiendo el trabajo en varias tareas. Estas
tareas pueden ser ejecutadas de manera secuencial o paralela.

En esta actividad se implementó en Python un contador de palabras inspirado
en MapReduce. El programa compara tres formas de ejecución: secuencial,
mediante hilos y mediante procesos.

\section{Objetivo}

Implementar un contador de palabras que represente las principales etapas
de MapReduce y utilizar las herramientas de concurrencia de Python para
comparar la ejecución secuencial, con hilos y con procesos.

\section{Descripción del programa}

El programa recibe un texto y cuenta cuántas veces aparece cada palabra.
Antes de realizar el conteo, convierte todas las palabras a minúsculas y
elimina algunas palabras comunes, como ``el'', ``la'', ``de'' y ``en''.

La información se procesa mediante las siguientes etapas:

\begin{enumerate}
    \item División del texto en varios fragmentos.
    \item Aplicación de la función Map sobre cada fragmento.
    \item Agrupación de los resultados mediante Shuffle.
    \item División de los grupos entre los reductores.
    \item Suma de las apariciones mediante Reduce.
    \item Unión de los resultados parciales.
\end{enumerate}

\section{Etapas de la implementación}

\subsection{División del texto}

La función \texttt{split\_text()} divide el texto utilizando sus líneas.
Cada fragmento puede ser enviado a un trabajador diferente.

La división por líneas evita cortar una palabra por la mitad y permite que
las tareas procesen partes independientes del archivo.

\subsection{Etapa Map}

La función \texttt{map\_chunk()} recibe un fragmento del texto. Mediante una
expresión regular identifica las palabras y genera el par
\texttt{(palabra, 1)} para cada una.

Por ejemplo, para la entrada:

\begin{center}
\texttt{Python procesa datos}
\end{center}

La función Map produce:

\begin{center}
\texttt{(python, 1), (procesa, 1), (datos, 1)}
\end{center}

\subsection{Etapa Shuffle}

La función \texttt{shuffle()} agrupa los valores que tienen la misma palabra.
Por ejemplo:

\begin{center}
\texttt{(python, 1), (python, 1)}
\end{center}

se transforma en:

\begin{center}
\texttt{(python, [1, 1])}
\end{center}

Esta etapa se realiza en el proceso principal antes de enviar los grupos a
los reductores.

\subsection{Partición}

La función \texttt{partition()} divide las palabras agrupadas en varias
partes. Cada parte puede ser procesada por un trabajador distinto durante
la etapa Reduce.

\subsection{Etapa Reduce}

La función \texttt{reduce\_partition()} suma la lista de valores de cada
palabra.

Por ejemplo:

\begin{center}
\texttt{(python, [1, 1])} $\longrightarrow$ \texttt{(python, 2)}
\end{center}

Finalmente, la función \texttt{merge()} combina todos los diccionarios
parciales en un único resultado.

\section{Modos de ejecución}

El programa utiliza tres modos de ejecución:

\begin{itemize}
    \item \textbf{Secuencial:} todas las funciones se ejecutan una después de
    otra en el proceso principal.

    \item \textbf{Hilos:} utiliza \texttt{ThreadPoolExecutor} para distribuir
    las tareas entre varios hilos del mismo proceso.

    \item \textbf{Procesos:} utiliza \texttt{ProcessPoolExecutor} para ejecutar
    las tareas en procesos separados.
\end{itemize}

Los procesos son útiles para trabajos que utilizan principalmente el
procesador. Los hilos suelen ser más apropiados cuando las tareas pasan
tiempo esperando operaciones de entrada y salida.

\section{Código fuente}

El siguiente código contiene la implementación completa del contador de
palabras:

\lstinputlisting[
    caption={Implementación de MapReduce en Python},
    label={lst:mapreduce}
]{\actividadactual/mapreduce.py}

\section{Ejemplo de ejecución}

El programa puede ejecutarse desde la terminal con el siguiente comando:

\begin{verbatim}
python3 mapreduce.py
\end{verbatim}

También se puede indicar un archivo de texto:

\begin{verbatim}
python3 mapreduce.py archivo.txt
\end{verbatim}

Una ejecución de prueba produjo el siguiente resultado:

\begin{verbatim}
Available cores used: 4
Text size: 403 characters

sequential    0.00004 s   (41 unique words)
threads       0.00032 s   (41 unique words)
processes     0.07331 s   (41 unique words)

Top 10 most frequent words:
  come            2
  programacion    2
  cada            1
  cardillo        1
  carne           1
  ciguena         1
  combina         1
  comia           1
  detras          1
  diccionario     1
\end{verbatim}

Los tiempos pueden cambiar dependiendo de la computadora y de los programas
que estén ejecutándose en ese momento.

\section{Análisis de resultados}

Los tres modos encontraron 41 palabras diferentes y produjeron exactamente
el mismo conteo. Esto demuestra que la división del trabajo no modifica el
resultado final.

En esta prueba la versión secuencial fue la más rápida porque el texto es
pequeño. La creación y coordinación de hilos y procesos necesita tiempo
adicional. Con archivos más grandes, la versión con procesos puede aprovechar
mejor los diferentes núcleos del procesador.

La versión con hilos no necesariamente mejora el tiempo en este ejercicio,
debido a que el conteo de palabras es una operación que utiliza principalmente
el procesador.

\section{Conclusión}

La actividad permitió implementar las etapas principales de MapReduce:
división, Map, Shuffle, partición y Reduce. También permitió comparar la
ejecución secuencial con el procesamiento mediante hilos y procesos.

El programa genera resultados correctos en los tres modos. Sin embargo, el
paralelismo no siempre significa una ejecución más rápida, especialmente
cuando la cantidad de información es pequeña y el costo de crear trabajadores
es mayor que el trabajo realizado.